% Prijava teme seminarske naloge -- NPMP 2026/2027, kvantno procesiranje
% Skupina 20. Prevajanje: pdflatex predlog_teme.tex
\documentclass[10pt,a4paper]{article}
\usepackage[utf8]{inputenc}
\usepackage[T1]{fontenc}
\usepackage[slovene]{babel}
\usepackage{mathptmx}
\usepackage{courier}
\usepackage[a4paper,margin=1.7cm,top=1.4cm,bottom=1.4cm]{geometry}
\usepackage{amsmath}
\usepackage{physics}
\usepackage{booktabs}
\usepackage{enumitem}
\usepackage[hidelinks]{hyperref}
\setlist{nosep,leftmargin=1.2em}
\setlength{\parindent}{0pt}
\setlength{\parskip}{3pt}
\pagestyle{empty}
\sloppy

\begin{document}

{\large\textbf{Prijava teme seminarske naloge -- NPMP 2026/27, kvantno procesiranje}}\\
\textbf{Skupina 20:} Dominik Krašovec, Erik Radovičevič \hfill Ljubljana, [DATUM ODDAJE]

\medskip
\textbf{Naslov:} \emph{Meja prepletenosti na procesorju IBM: koliko kubitov lahko zares prepletemo v stanje GHZ?}

\textbf{Raziskovalno vprašanje.} Koliko kubitov lahko procesor IBM Heron (ibm\_fez, 156 kubitov) pripravi v stanje GHZ
$\ket{\mathrm{GHZ}_N}=\tfrac{1}{\sqrt2}(\ket{0}^{\otimes N}+\ket{1}^{\otimes N})$ z zvestobo $F>0{,}5$
(zadosten pogoj za \emph{pravo} večdelno prepletenost vseh $N$ kubitov) in za koliko kubitov to mejo $N_{\max}$
premakneta (a)~zasnova vezja (linearna veriga proti drevesu z majhno globino) ter (b)~dinamično razklapljanje (DD) in ublažitev napak branja?

\textbf{Hipoteza.} Hipoteza je napoved klasičnega modela proračuna napak (vrata, čakanje, branje) iz kalibracije FakeFez, ki se s
šumno simulacijo pri $N\le13$ ujema na največ 0,016. Največji $N$ z napovedanim $F>0{,}5$: linearna veriga $N_{\max}\approx 26$
brez ublažitve in $\approx 33$ z ublažitvijo branja, drevo z globino $\propto\!\sqrt N$ pa $\approx 33$ oziroma $\approx 70$. Surovo mejo določajo predvsem napake branja, mejo verige pa
čakanje kubitov. DD bolj pomaga verigi, ker njeni kubiti dlje čakajo; šumni model pozna le markovski T1/T2, zato učinek DD
izmerimo samo na procesorju.

\textbf{Metrike.}
\begin{itemize}
\item Zvestoba $F=(P+C)/2$: populacija $P=p(0\dots0)+p(1\dots1)$, koherenca $C=2\abs{c_N}$, kjer je $c_N$ Fourierova komponenta
      frekvence $N$ paritete $\expval{\Pi(\phi)}$ po rotaciji vseh kubitov v bazo $\cos\phi\,X+\sin\phi\,Y$, izmerjene pri 9 fazah
      $\phi_k=\pi k/9$. Pariteta vsebuje le frekvence $N, N-2,\dots$, zato diskretna Fourierova transformacija (DFT) loči signal od
      šumnih členov (ostanejo le členi z 9, 18, \dots{} hkratnimi preklopi; $N$ ni večkratnik 9).
\item $F$ surovo in z lokalno ublažitvijo branja (matrika $2\times2$ na kubit iz kalibracijskih vezij $\ket{0\dots0}$ in
      $\ket{1\dots1}$, po 4000 posnetkov). Negotovost z zankanjem (\emph{bootstrap}), ki na novo vzorči tudi kalibracijo.
\item $N_{\max}$: največji $N$ s $F-2\sigma_F>0{,}5$, za vsako od 4 konfiguracij (veriga/drevo $\times$ brez/z DD), surovo in
      ublaženo; med izmerjenimi $N$ ga podamo kot interval. Prirastek $\Delta N_{\max}$ zaradi zasnove, DD in ublažitve.
\item Lastnosti vezja pred prevajanjem in po njem (globina, število vrat CZ), po prevajanju še trajanje; razlika med modelom in procesorjem.
\end{itemize}

\textbf{Načrt poskusov.}
\begin{enumerate}
\item \emph{Vezja:} veriga po najcenejši najdeni poti v grafu povezav (iskanje v globino; cena je vsota napak branja in vrat iz kalibracije), drevo z iskanjem v
      širino (BFS) iz korena z najmanjšo oceno napak in globine. Obe brez dodanih vrat SWAP ($N-1$ vrat CZ).
\item \emph{Idealni simulator} (Qiskit Aer, metoda MPS): preverba $F=1$. \emph{Šumni simulator} FakeFez (kalibracija
      26.\,2.\,2025) do $N\approx 13$--$20$ in model proračuna napak za vse $N$. Ponovitev prejšnjega dela: naš 5-kubitni zagon
      na ibm\_fez ($P=0{,}855$) proti FakeFez.
\item \emph{Pilot} na ibm\_fez (1 posel, 88 vezij): veriga $N=10, 22$, drevo $N=20, 40$, vsaka točka brez in z DD
      (impulza X v času čakanja); točka ima 1 populacijsko in 9 paritetnih vezij, posel še 8 kalibracijskih. Preveri model porabe.
\item \emph{Glavni zagon} (1 posel, 148 vezij): veriga $N=16, 28, 34$, drevo $N=30, 46, 50, 56$ (drevo po pilotu premaknjeno z
      $N=30, 50, 64, 76$). Para brez/z DD sta v istem poslu, zato ju ne loči sprememba kalibracije.
\item Zgostitev točk okoli $N_{\max}$ (1 posel, 88 vezij) in dodatni posel brez kubitov s poslabšanim branjem (68 vezij). Primerjava: simulator / model / procesor ter objave
      (18 kubitov, $F=0{,}517$, Wei idr.\ 2020; 27 kubitov, $F=0{,}546$, Mooney idr.\ 2021; 120 kubitov, $F=0{,}56$, IBM 2025).
\end{enumerate}

\textbf{Ocena procesorskega časa.} Izmerjeno 2.\,11.\,2026: 5-kubitni GHZ, 1 vezje $\times$ 1000 posnetkov na ibm\_fez
(job ID \texttt{dbk48fl6p51s73clupeg}) je porabil \textbf{2~s} (ocena IBM 2,27~s). To se ujema z modelom IBM
$t \approx 2\,\mathrm{s}\cdot(\text{št.\ poslov}) + 0{,}35\,\mathrm{ms}\cdot(\text{št.\ vezij}\times\text{št.\ posnetkov})$,
zato vezja združujemo v malo poslov (1000 posnetkov na vezje).
\begin{center}\small
\begin{tabular}{lrrr}
\toprule
Del poskusa & vezja & posli & zaračunano \\
\midrule
prvi preizkus (2.\,11.\,2026) & 1 & 1 & 2~s \\
pilot & 88 & 1 & 25~s \\
glavni zagon & 148 & 1 & 41~s \\
zgostitev okoli $N_{\max}$ & 88 & 1 & 25~s \\
dodatni posel & 68 & 1 & 20~s \\
\midrule
\textbf{skupaj} & \textbf{393} & \textbf{5} & \textbf{113~s} \\
\bottomrule
\end{tabular}
\end{center}
Na voljo je bilo 600~s (Open Plan: drseče 28-dnevno okno). Vsak posel ima zgornjo mejo trajanja (\texttt{max\_execution\_time}), ki jo
skripta pred oddajo rezervira v dogovorjenem proračunu.

\textbf{Orodja.} Python~3, Qiskit~2.6 (vezja, prevajanje, razporejanje ALAP, DD), Qiskit Aer 0.18 (MPS, šumni model FakeFez),
qiskit-ibm-runtime 0.51 (SamplerV2), NumPy, Matplotlib, pytest; IBM Quantum Open Plan; \LaTeX{} (predloga zbornika); GitHub.\par

\textbf{Razdelitev dela.} Oba člana (Dominik in Erik) delava na vseh delih naloge: gradnja vezij (veriga, BFS drevo) in izbira
kubitov, simulacije (idealni, FakeFez) in model napak, zagoni na IBM in beleženje job ID ter porabe časa, analiza ($P$, $C$, $F$,
\emph{bootstrap}, ublažitev branja), pregled literature, uvod, razprava, prosojnice in repozitorij (README). Tako oba razumeva celotno delo.
Vsak del napiše ali izvede eden, drugi ga pregleda in neodvisno ponovi, vlogi pa si po delih izmenjava.

\end{document}
